\documentclass[11pt]{article}

% --- packages (all standard; compiles with pdflatex / Overleaf) ---
\usepackage[utf8]{inputenc}
\usepackage[T1]{fontenc}
\usepackage[a4paper,margin=2cm]{geometry}
\usepackage{booktabs}
\usepackage{xcolor}
\usepackage{titlesec}
\usepackage{enumitem}
\usepackage[hidelinks]{hyperref}

\definecolor{metroRed}{HTML}{E4002B}
\definecolor{metroGreen}{HTML}{00A651}
\definecolor{ink}{HTML}{0B1B3A}
\definecolor{gold}{HTML}{B8860B}

\titleformat{\section}{\large\bfseries\color{metroRed}}{\thesection}{0.6em}{}
\titleformat{\subsection}{\normalsize\bfseries\color{ink}}{\thesubsection}{0.6em}{}
\setlist{nosep,leftmargin=1.4em,topsep=2pt}
\renewcommand{\arraystretch}{1.2}

\newcommand{\q}[1]{\medskip\noindent\textbf{Q. #1}\par\nopagebreak}
\newcommand{\an}[1]{\noindent\textit{A.} #1\par}
\newcommand{\term}[1]{\noindent\textbf{\color{ink}#1} --- }
\newcommand{\elem}[1]{\smallskip\noindent\textbf{\color{metroRed}#1}\ \ }

\begin{document}

\begin{center}
  {\Huge\bfseries\color{metroRed} Dubai Metro Forecasting}\\[3pt]
  {\LARGE\bfseries Complete Presentation Guide \& Q\,\&\,A}\\[6pt]
  {\normalsize \textbf{Team:} Kartik \;\textbullet\; Prem \;\textbullet\; Gagandeep \;\textbullet\; Samuel}\\[3pt]
  {\small Every screen, button, number and concept on the dashboard --- explained so simply that a
  high-school student can read it once and present confidently.}\\[4pt]
  {\color{metroGreen}\rule{0.6\textwidth}{1.2pt}}
\end{center}

% ============================================================
\section{Who presents what}
We present \emph{using the live dashboard}. It has four tabs across the top, and each teammate owns one.

\begin{center}
\begin{tabular}{llp{8.2cm}}
\toprule
\textbf{Order} & \textbf{Presenter} & \textbf{Owns}\\
\midrule
1 & \textbf{Kartik}    & Opening \& setting the stage + the \textbf{Live} (landing) page.\\
2 & \textbf{Prem}      & The \textbf{Stations} page.\\
3 & \textbf{Gagandeep} & The \textbf{Network} page (incl.\ the live \textbf{Crowd Alert}).\\
4 & \textbf{Samuel}    & The \textbf{Model} page (algorithms \& results).\\
\bottomrule
\end{tabular}
\end{center}

\noindent\textbf{Three golden rules:} (1) start each part with one plain sentence of what the
viewer is looking at; (2) the clock is the \emph{real} Dubai time --- if it is night, use the
\textbf{Demo} button to jump to a busy hour; (3) never bluff a number --- every number you need is on
the screen or in this guide.

% ============================================================
\section{The bar across the top (always on screen) \textnormal{--- Kartik introduces}}
This top bar (the ``navbar'') stays on every tab.
\begin{itemize}
  \item \elem{Logo / title} ``Dubai Metro --- Next-Hour Forecasting --- LSTM + RNN.'' Tells everyone
        what the tool is and the model behind it.
  \item \elem{Four tabs} \textbf{Live, Stations, Network, Model} --- the four parts of our talk. Click
        to switch pages.
  \item \elem{Demo button} (labelled ``Demo'') --- a \emph{presentation aid}. Click it to open a small
        panel where we can pretend it is any hour of the day (and weekday or weekend), so we can show
        the morning rush even if it is currently night. \textbf{Off by default} --- normally the app
        uses the real time. When on, the clock says ``DEMO --- simulated'' so we are being honest.
  \item \elem{Live clock} shows the current \textbf{Dubai time} (GST, which is UTC$+$4, no daylight
        saving). A \textbf{green dot} means trains are running (``GST --- live''); a \textbf{grey dot}
        means ``service closed'' (night).
\end{itemize}

\noindent\textbf{Mention once (global safety feature):} a \textbf{Crowd Alert} can pop up at the top
of \emph{any} tab whenever the hour changes and some stations are about to get busy. Gagandeep will
explain it in detail; just know it can appear at any time.

% ============================================================
\section{Part 1 --- Live page \textnormal{(Kartik)}}
\textbf{Opening line:} ``Our project predicts how many people will enter each Dubai Metro station in
the \emph{next hour}, live, so operators can add trains \emph{before} a crowd forms instead of
reacting after it.''

\textbf{The story \& data (say this first):}
\begin{itemize}
  \item \textbf{Why it matters:} too few trains $\rightarrow$ dangerous crowding and long waits; too
        many $\rightarrow$ wasted energy and money. Knowing demand one hour ahead lets the operator
        get it right.
  \item \textbf{Real data:} Dubai's official open data (RTA, via the ``Dubai Pulse'' portal) ---
        real smart-card tap records. After cleaning: \textbf{about 503{,}000 entries, 42 days,
        53 stations, all from 2026} (the latest).
  \item \textbf{Proof it's real:} the busiest stations the data finds --- BurJuman, Union, Al~Rigga,
        Mall of the Emirates, Burj~Khalifa/Dubai~Mall --- are Dubai Metro's actual busy stations.
\end{itemize}

\subsection{Every element on the Live page}
\elem{Status badge (top-left)} ``LIVE --- NEXT-HOUR FORECAST'' when trains run; ``METRO CLOSED'' at
night; ``Service winding down'' in the last operating hour.

\elem{Big clock} the live time, ticking by the second, with the weekday and ``Dubai GST''.

\elem{Headline label} ``Forecast network inflow for HH:00'' --- the hour we are predicting (always the
\emph{next} hour).

\elem{The big number} the predicted total check-ins across all 53 stations for that next hour. Its
\textbf{colour} is a quick demand level: green = Low, blue = Moderate, amber = High, red = Peak.

\elem{Sub-line} ``check-ins across all 53 stations --- [Low/Moderate/High/Peak] demand --- up/down X\%
vs this hour.'' The percentage says whether the next hour is busier or quieter than the current one.

\elem{Three pills} (small tags): ``This hour (HH:00): N'' (current hour's figure), ``Model:
LSTM+RNN'' (the model in use), and ``R\textsuperscript{2} 87\%'' (its accuracy --- explained below).

\elem{Right-hand chart --- ``Today's expected demand --- network''} two lines over the day:
\begin{itemize}
  \item \textbf{Typical} (shaded blue) --- the \emph{usual} demand for this kind of day (the average
        from history). Think ``what normally happens.''
  \item \textbf{Forecast} (red line) --- what our model predicts. It sits right on top of the typical
        curve, which shows the model is accurate.
  \item \textbf{Green dot ``now''} --- marks the current hour, so you can see where we are in the day.
  \item Caption tells you whether it is using the \textbf{Weekday} or \textbf{Weekend} profile.
\end{itemize}

\elem{Four KPI cards (bottom strip)} --- the project's headline scoreboard:
\begin{enumerate}
  \item \textbf{Forecast accuracy (R\textsuperscript{2}) $\approx$87\%} --- of all the up-and-down
        movement in real demand, our model explains about 87\% of it. Higher is better (100\% is perfect).
  \item \textbf{Network tracking (corr) $\approx$0.83} --- on a real test day, how tightly the forecast
        line follows the actual line (1.0 = move in perfect lock-step).
  \item \textbf{Stations forecast $=$53} --- we cover \emph{every} station, not just a few famous ones.
  \item \textbf{Real taps --- 2026 $\approx$503k} --- the size of the real dataset we trained on.
\end{enumerate}

\elem{Closed state (if it is night)} the big number becomes ``-- --'' and it says ``Dubai Metro is
not running now. First trains at 05:00.'' \emph{This is correct behaviour} --- use the Demo button to
jump to a daytime hour for the demo.

% ============================================================
\section{Part 2 --- Stations page \textnormal{(Prem)}}
\textbf{Opening line:} ``The Live page gave the whole-network picture. This page lets us drill into
\emph{any single station} --- all 53 of them.''

\subsection{Every element on the Stations page}
\elem{Search box} ``Search any station...'' --- type part of a name (e.g.\ ``max'', ``Energy'',
``Mall'') to filter the list instantly.

\elem{Station list (left)} every station as a row with: a \textbf{coloured dot} (red = Red Line,
green = Green Line), the \textbf{station name}, and its \textbf{next-hour forecast number} on the
right (shows ``--'' when the metro is closed). Click a row to select it.

\elem{Detail panel (right) --- header} the chosen station's name and a \textbf{line tag} (``Red'' or
``Green''), plus a caption: ``Next-hour forecast for HH:00 --- anchored to live Dubai time HH:MM''
(or a closed-message at night).

\elem{Big number + label} that station's predicted entries for the next hour, with a demand word
(Peak / High / Moderate / Low).

\elem{Detail chart} the same two-line idea as the Live page but for this one station: \textbf{Typical}
demand (blue) vs our \textbf{Forecast} (red), with the green \textbf{now} marker. You can literally
watch the model trace the station's daily rhythm.

\subsection{Concepts to mention here}
\term{Station embedding} the model learns a little ``ID card'' of numbers for each station, so it
treats a giant interchange (BurJuman) differently from a quiet stop --- that is how \emph{one} model
serves all 53 stations well instead of needing 53 separate models.

\term{Line colours} Dubai Metro has a \textbf{Red Line} and a \textbf{Green Line}; we keep that colour
coding everywhere so the audience instantly knows which line a station is on.

\textbf{Nice thing to show:} pick a small station and a big interchange and compare their curves ---
same model, very different shapes, both tracked well.

% ============================================================
\section{Part 3 --- Network page \textnormal{(Gagandeep)}}
\textbf{Opening line:} ``Operators don't watch one station --- they watch the whole network. This page
turns forecasts into action, and it's where our live \emph{Crowd Alert} lives.''

\subsection{Every element on the Network page}
\elem{``Crowding outlook --- every station''} the top card. For the \emph{upcoming} hour it ranks all
53 stations busiest-first, each with:
\begin{itemize}
  \item a coloured dot (line), the station name, a \textbf{coloured bar} showing how full it is vs its
        own busy peak, the \textbf{forecast number}, and a \textbf{tier word}: Low (green),
        Moderate (blue), High (amber), Peak (red).
  \item a \textbf{search box} to find any station;
  \item the caption notes the comparison uses a \textbf{real train capacity}: \textbf{643 passengers
        per 5-car train} (RTA's official figure).
\end{itemize}

\elem{``Real Dubai demand --- average hourly inflow''} a chart with two shaded curves: \textbf{Weekday}
(red) and \textbf{Weekend} (green) average demand by hour. It clearly shows the \textbf{morning peak
(07--09)} and \textbf{evening peak (16--19)}, and that weekends look different.

\elem{``Busiest stations''} a ranked bar list (top 8) by real check-ins, matching Dubai Metro's known
hotspots --- more proof the data is genuine.

\subsection{The Crowd Alert (our live safety feature)}
\elem{What it is} a pop-up that appears at the top of \emph{any} tab the moment the hour changes and
some stations are forecast to get busy. It says, e.g., ``\emph{Crowd alert --- 08:00: High demand
expected at Al~Rigga, Union, Abu~Baker Al~Siddique --- consider adding trains / staff},'' with a small
chip per station showing the expected number.
\elem{How to demo it} open \textbf{Demo}, turn it on, and drag the hour to \textbf{07} or \textbf{16};
when the hour ticks to the rush hour, the alert fires. It auto-closes after about 10 seconds, or click
the ``x''. It stays hidden when the metro is closed.
\elem{Why it matters} this is the project's whole point made visible: the system warns staff
\emph{before} the surge, on whatever screen they happen to be looking at.

\subsection{Concepts to mention here}
\term{Crowding ratio} forecast demand divided by a station's own busy level --- a simple ``how full
will it be'' score that turns a raw number into Low/High/Peak.
\term{The Dubai ``arrives-full'' story} a train coming from the far (Expo) end can already be packed,
so even a smaller station can be hard to board at peak --- which is exactly why we forecast
\emph{every} station, not just the famous ones.

% ============================================================
\section{Part 4 --- Model page \textnormal{(Samuel)}}
\textbf{Opening line:} ``Here is the engine and the proof. We compared many algorithms \emph{fairly};
our LSTM+RNN wins, and it clearly beats the simple baselines.''

\subsection{Every element on the Model page}
\elem{``Approach comparison'' table} every model on the same test data (about 7{,}000 held-out
hours). Columns are the four accuracy scores (MAE, RMSE, MAPE, R\textsuperscript{2}). Our
\textbf{LSTM+RNN} row is highlighted and tagged \textbf{``OUR SOLUTION''}; the two simplest methods are
tagged \textbf{``baseline''}. Caption reminds: \emph{lower error = better}.

\elem{``Validation --- real day, actual vs forecast'' chart} on a real day the model never saw, the
\textbf{Actual} line (blue) vs our \textbf{Forecast} (red dashed). They track closely
(\textbf{correlation $\approx$0.83}). The caption stresses this is on \emph{unseen real} data --- the
honest test.

\elem{``How the model works'' (4 steps)} a plain-English recipe: (1) real 2026 data $\rightarrow$
hourly counts; (2) features (station ID, hour, weekend, recent counts, the ``usual'' value);
(3) the LSTM+RNN hybrid reads the last 6 hours and predicts the next; (4) result: best scores,
beats the baselines.

\subsection{The headline numbers (memorise)}
\begin{center}
\begin{tabular}{lcccc}
\toprule
\textbf{Approach} & \textbf{MAE} & \textbf{RMSE} & \textbf{MAPE} & \textbf{R\textsuperscript{2}}\\
\midrule
\textbf{\color{metroRed}LSTM+RNN (ours)} & \textbf{4.45} & \textbf{8.09} & \textbf{30\%} & \textbf{0.87}\\
CNN-LSTM / RNN / LSTM / GRU & 4.5--4.7 & 8.1--8.7 & $\sim$31\% & 0.85--0.87\\
Naive (persistence)        & 4.99 & 8.88 & 38\% & 0.85\\
Climatology                & 12.5 & 23.6 & 75\% & below 0\\
\bottomrule
\end{tabular}
\end{center}
\noindent\textbf{One-liner:} ``Lower is better. Ours has the lowest error, and is much better on
percentage error than the simple `repeat the last hour' guess.''

% ============================================================
\section{Plain-English dictionary (every concept)}
\term{Forecasting} guessing what comes next from what already happened (like guessing tomorrow's
weather from this week's).
\term{Passenger inflow} how many people \emph{tap in} (enter) at a station in a time slot.
\term{Time series} a list of numbers in time order (people each hour, hour after hour).
\term{AFC (Automated Fare Collection)} the smart-card tap system; every tap is one timestamped record.
\term{Neural network} a computer model that learns patterns from many examples instead of fixed rules.
\term{RNN (Recurrent Neural Network)} a neural network \emph{with memory}; it reads a sequence one
step at a time and remembers recent steps --- like reading a sentence word by word. Perfect for time series.
\term{LSTM (Long Short-Term Memory)} a smarter RNN with little ``gates'' that choose what to
\textbf{keep}, \textbf{forget} and \textbf{output}, so it remembers important things for longer ---
like a student with a good notebook who keeps key facts and drops the noise.
\term{GRU (Gated Recurrent Unit)} a lighter, faster cousin of the LSTM with fewer gates; similar idea.
\term{CNN (Convolutional Neural Network)} usually for images; slides a small window to find \emph{local}
patterns. In \textbf{CNN-LSTM}, the CNN spots short bursts and the LSTM tracks them over time.
\term{Our hybrid (LSTM + RNN)} we run an LSTM \emph{and} a SimpleRNN on the same input side-by-side and
join their answers --- LSTM for longer memory, RNN for a quick read of recent momentum.
\term{Ensemble} ask several copies of the model and average their answers --- like averaging several
friends' guesses. Averaging cancels random luck, so results are steadier and a bit better. Ours
averages three.
\term{Station embedding} a small learned ``ID card'' of numbers per station, so one model can
specialise for each of the 53 stations.
\term{Climatology} the ``usual'' pattern --- the historical average for that station, hour and
weekday. We feed it as a hint and also use it as a baseline.
\term{Bias calibration} a final nudge so forecasts aren't always slightly too low or too high ---
like fixing a ruler that always reads 1\,cm short.
\term{Features} the clues we give the model: hour, day-of-week, weekend flag, recent counts, station
ID and the ``usual'' value.
\term{Training / Validation / Test} study notes / practice quiz / final exam. We score accuracy on
the \emph{test} days the model never saw, so the score is honest.
\term{Data leakage} accidentally letting the model peek at exam answers while studying. We avoided it:
the ``usual'' values and calibration come only from training/validation days, and we split by date.
\term{Overfitting} memorising the training data instead of learning the real pattern (does great in
practice, badly on the exam). The ensemble, dropout and an honest test guard against it.

% ============================================================
\section{Every algorithm in the comparison (what \& why)}
\begin{itemize}
  \item \textbf{Naive (persistence)} --- the simplest guess: ``next hour = this hour.'' The bar every
        real model must beat.
  \item \textbf{Climatology} --- ``next hour = the usual for this hour/day.'' Captures the typical
        shape but ignores how busy \emph{today} is, so it's weak alone (even scores negative
        R\textsuperscript{2}).
  \item \textbf{RNN (only)} --- a plain recurrent network reading the last 6 hours.
  \item \textbf{LSTM (only)} --- a stacked LSTM; better long memory than a plain RNN.
  \item \textbf{GRU} --- a lighter gated network; fast and competitive.
  \item \textbf{CNN-LSTM} --- a CNN finds short patterns, an LSTM tracks them over time.
  \item \textbf{LSTM+RNN (ours)} --- the parallel hybrid, deployed as a 3-model ensemble with bias
        calibration. Lowest error overall --- the model the dashboard uses.
\end{itemize}
\noindent All learning models share the \emph{same} inputs; only the sequence core differs --- so the
comparison is fair.

% ============================================================
\section{The four accuracy scores (MAE, RMSE, MAPE, R\textsuperscript{2})}
An ``error'' = predicted minus actual passengers for one hour at one station.
\term{MAE (Mean Absolute Error)} the average error size, ignoring sign. ``On average we're off by
about 4--5 passengers per station-hour.'' \textbf{Lower is better.}
\term{RMSE (Root Mean Squared Error)} like MAE but \emph{punishes big misses more}. If it's much larger
than MAE, a few large mistakes exist. \textbf{Lower is better.}
\term{MAPE (Mean Absolute Percentage Error)} the error as a \emph{percentage} of the real value
(``about 30\% off''). Fair for comparing busy and quiet stations. \textbf{Lower is better.}
\term{R\textsuperscript{2} (R-squared)} how much of the up-and-down pattern the model explains, from
0 to 1. $1$ = perfect; $0$ = no better than always guessing the average; \emph{negative} = worse than
that. Ours $\approx$0.87 means we explain $\approx$87\%. \textbf{Higher is better.}

\medskip\noindent\textit{If asked ``MAE is close to others --- why do you win?''} our \textbf{MAPE} is
clearly better, meaning we're much more accurate at the quieter stations and off-peak hours where
``just repeat the last hour'' fails; and we lead on \textbf{RMSE}, which matters most for the big,
safety-critical peaks.

% ============================================================
\section{Likely questions --- easy to hard (with answers)}
Say the answers in your own words; short is fine.

\subsection{Easy (definitions \& basics)}
\q{In one line, what does your project do?}
\an{It predicts how many passengers will enter each Dubai Metro station in the next hour, live.}
\q{What is forecasting?}
\an{Using past data to estimate a future value --- here, next hour's entries.}
\q{Is your data real?}
\an{Yes --- official RTA smart-card tap records from Dubai's open-data portal (Dubai Pulse), 2026.}
\q{What is an RNN?}
\an{A neural network with memory that reads a sequence step by step --- ideal for time-ordered data.}
\q{What is an LSTM?}
\an{A special RNN with gates that choose what to remember, forget and output, so it handles long
patterns better than a plain RNN.}
\q{What do MAE and R\textsuperscript{2} mean?}
\an{MAE = average error size (lower better, ours $\approx$4.5 passengers). R\textsuperscript{2} =
fraction of the pattern explained (higher better, ours $\approx$0.87, i.e.\ 87\%).}
\q{What are the four KPI cards on the landing page?}
\an{Accuracy (R\textsuperscript{2} $\approx$87\%), Network tracking (correlation $\approx$0.83),
Stations forecast (53), and Real taps in 2026 ($\approx$503k).}
\q{What does the Crowd Alert do?}
\an{When the hour changes, it pops up on any tab and names the stations about to get crowded, so staff
can act early.}
\q{Why Dubai?}
\an{A modern automated metro with rich tap data, event/tourism-driven surges, and a smart-city agenda
where crowd forecasting fits perfectly.}

\subsection{Medium (design choices)}
\q{Why ``next hour'' and not next day?}
\an{Next-hour is the operational window --- it's what lets you add a train or open gates in time.}
\q{Why combine LSTM and RNN instead of one?}
\an{LSTM gives longer memory, the SimpleRNN a compact read of recent momentum; joining them and
averaging an ensemble is steadier and scores best in our fair comparison.}
\q{Difference between LSTM and GRU?}
\an{Both are gated RNNs. LSTM has three gates and a separate memory cell; GRU merges them into two ---
simpler and faster, often similar accuracy.}
\q{What is the naive baseline and why include it?}
\an{``Next hour = this hour.'' If a model can't beat it, the model is pointless. Ours beats it on every metric.}
\q{What is climatology and why is it a baseline?}
\an{The historical average for that station/hour/weekday --- the ``usual.'' It misses how busy today
is, so alone it's weak (negative R\textsuperscript{2} on our test).}
\q{Why use only 2026 data?}
\an{About 90\% of available days are 2026 at today's ridership; the few old days are much lower-volume.
Mixing them creates a train/test mismatch and teaches outdated patterns.}
\q{What is an ensemble and why does averaging help?}
\an{Each trained copy differs slightly due to randomness; averaging cancels that noise, lowering error
and making the ranking stable run-to-run.}
\q{How is the live clock built, and what happens at night?}
\an{We compute Dubai time directly (UTC$+$4, no daylight saving). At night it shows ``closed, first
trains 05:00''; in service hours it shows the next-hour forecast. A demo override simulates any hour
for presentations.}
\q{How did you avoid data leakage?}
\an{We split by date (train on earlier days, test on later days) and compute climatology and
calibration only from training/validation data --- never from the test days.}

\subsection{Advanced (technical depth)}
\q{Your forecast first looked $\sim$15\% too low. Why, and how did you fix it?}
\an{We train with the Huber loss, which targets the \emph{median}. Hourly inflow is right-skewed (many
small values, occasional big peaks), so the median sits below the mean and the model under-predicts.
We fixed it with a multiplicative bias-calibration factor (separate for weekday/weekend) fit only on
the validation set, making forecasts unbiased without touching test data.}
\q{Why is climatology's R\textsuperscript{2} negative?}
\an{Each sampled day has a different overall level; a fixed historical average can be systematically
off for a given day, so it does worse than simply predicting the global mean --- which is what a
negative R\textsuperscript{2} means. The recurrent model fixes this by reading how busy today is.}
\q{Why a chronological split, not a random one?}
\an{It's a time-series task --- in real life you predict the future from the past. A random split could
let the model see future days while training, inflating the score. Chronological keeps it honest.}
\q{Why does one-step next-hour forecasting lag at the very tip of a peak?}
\an{It predicts each hour from the hours just before; right at a sharp peak those are still rising, so a
single step slightly under-shoots the exact peak. Our climatology hint and ``today's level'' feature
reduce it (network tracking $\approx$0.83; the typical curve overlays at $\approx$0.99). Multi-step
sequence-to-sequence / attention models predict the whole horizon at once --- our future work.}
\q{What exactly is the station embedding?}
\an{A lookup table mapping each station ID to a small learned vector (8 numbers). Training shapes these
so similar stations get similar vectors --- letting one shared model specialise per station.}
\q{Stations affect each other --- why not a Graph Neural Network (DCRNN/STGCN)?}
\an{Spatial graph models capture how a surge spreads between connected stations and are a strong next
step. We deliberately built a robust, explainable recurrent model first; the graph version is in our
future-work list.}
\q{What are the limitations of your data?}
\an{The open extract is a \emph{sample} of taps, so absolute counts are scaled and a day's level
varies. We handle it with per-station scaling and calibration, and we're transparent that shapes and
rankings are real even if raw magnitudes are sampled. Full data would sharpen the absolute numbers.}
\q{How would you take this to production?}
\an{Stream live AFC taps into the same hourly aggregation, run the trained ensemble each hour, and push
next-hour forecasts and crowd alerts to the control room; retrain periodically to track drift (new
lines, events, behaviour change).}
\q{Why parallel LSTM + RNN rather than stacking them?}
\an{Stacking a SimpleRNN after an LSTM sometimes discards useful detail and was unstable in our tests.
Running them in parallel and joining keeps both representations and reliably matches or beats either
alone --- especially as an averaged ensemble.}
\q{Which metric matters most here?}
\an{For safety/operations, \textbf{RMSE} (punishes big misses at busy times) and \textbf{MAPE}
(fairness across station sizes). Our model leads on both --- that's the real argument, not a tiny MAE edge.}

\bigskip
\begin{center}\small\color{ink}\rule{0.4\textwidth}{0.8pt}\\[4pt]
\textit{Speak simply, point at the screen, let the live demo and the Model-tab numbers do the arguing.}
\end{center}

\end{document}
